% --- SECCIÓ 5: CONCLUSIONS ---
\apartat{Conclusions}
Aquest treball ha permès aprofundir en la relació intrínseca entre la percepció sensorial i la capacitat comunicativa humana. Les conclusions principals s'estructuren segons els resultats obtinguts:

\begin{itemize}
    \item \textbf{Facilitació del llenguatge:} S'ha confirmat que els sistemes de comunicació augmentativa (CAA) actuen com un pont que potencia, en lloc de frenar, el desenvolupament de la intenció comunicativa.
    \item \textbf{Regulació sensorial:} Els estímuls controlats de la taula sensorial són eficaços per focalitzar l'atenció de l'usuari i reduir l'estrès ambiental, creant un estat òptim per a l'aprenentatge.
    \item \textbf{Autonomia i Empoderament:} La interacció directa amb la taula millora l'autoconcepte de l'usuari en veure que les seves accions tenen una resposta sensorial immediata i predictible.
\end{itemize}

En definitiva, la comunicació a través dels sentits és una eina clau per garantir la plena inclusió i dignitat de totes les persones en la nostra societat.
